%!TEX root = ../manuscript.tex

\begin{definition}[Nesting forest]\label{def:nesting-forest}
Let $w$ be a Gauss word with crossing labels $1,\dots,n$.
For two distinct crossings $a,b$, we say that $b$ strictly nests $a$
if the two positions of $a$ lie strictly between the two positions
of $b$ in the cyclic order.
We further say that $b$ \textbf{directly nests} $a$, written $a\prec b$,
if $b$ strictly nests $a$ and no third crossing $c$ satisfies
$a\prec c\prec b$.
The \textbf{nesting forest} $N(w)$ has vertex set $\{1,\dots,n\}$
and a directed edge $b\to a$ whenever $a\prec b$.
\end{definition}

\begin{lemma}[Partial order]\label{lem:D}
The strict nesting relation restricted to any subset of chords is a
partial order (antisymmetric and transitive).  Consequently the
nesting forest $N(w)$ is the Hasse diagram of this partial order
restricted to cover relations.
\end{lemma}

\begin{proof}
Antisymmetry: if $a$ is strictly nested in $b$ then the first position
of $b$ precedes the first position of $a$, so $b$ cannot simultaneously
be strictly nested in $a$.  Transitivity: if $a$ lies between the two
endpoints of $b$ and $b$ lies between the two endpoints of $c$, then
$a$ lies between the two endpoints of $c$.  Thus nesting is a strict
partial order on the chords of $w$.

The Hasse diagram of a finite partial order consists of the cover
relations $x\lessdot y$.  By definition, $a\prec b$ is exactly the
cover relation in the nesting order, so the directed edges of $N(w)$
are precisely the Hasse-diagram edges.
\end{proof}

\begin{lemma}[Unique parent]\label{lem:E}
Every chord that is strictly nested by at least one other chord has
exactly one direct parent in $N(w)$.
\end{lemma}

\begin{proof}
Let $a$ be a chord nested by some other chord; let $S$ be the set of
all chords that strictly nest $a$.  $S$ is non-empty and finite, and
the nesting relation totally orders $S$ (all chords in $S$ contain $a$
inside them, so they are pairwise nested).  Therefore $S$ has a unique
minimal element $b$ in the nesting order: the chord whose two
endpoints are closest to, but still surrounding, the endpoints of $a$.
If there were a chord $c$ with $a\prec c\prec b$, then $c$ would nest
$a$ but be strictly inside $b$, contradicting the minimality of $b$.
Hence $b$ directly nests $a$, i.e.  $a\prec b$, and $b$ is the unique
parent of $a$ in $N(w)$.
\end{proof}
